\section*{Entorno local}
\phantomsection
\addcontentsline{toc}{section}{Entorno local}
\label{sec:entorno-local}

El desarrollo de software contemporáneo no se apoya en un único entorno de
ejecución, sino en una sucesión de ellos que cada cambio recorre antes de llegar
al usuario final: la máquina del desarrollador, los entornos de verificación
automatizada, los de pruebas o preproducción y, finalmente, el de producción.
Cada uno de estos escalones cumple una función distinta y aporta un nivel
adicional de confianza sobre el cambio, pero todos comparten un mismo punto de
partida, ya que ningún cambio llega a los siguientes entornos sin haberse
escrito y ejecutado primero en el equipo de quien lo desarrolla. Esa posición
inicial convierte al entorno local en el lugar más económico para detectar un
error: cuanto más se parezca a producción, antes se manifiestan las
discrepancias que de otro modo aparecerían en etapas donde su diagnóstico y
corrección resultan considerablemente más costosos. En este sentido, Humble y
Farley \cite{Humble2010} sostienen que resulta esencial emplear el mismo proceso
de despliegue en todos los entornos —incluida la estación de trabajo del
desarrollador—, pues solo después de haberlo ejecutado centenares de veces sobre
entornos diversos cabe descartarlo como origen de un fallo. Los autores observan,
además, que la frecuencia con la que se despliega en cada entorno es inversa al
riesgo que este comporta: el desarrollador despliega continuamente, mientras que
el entorno en el que menos se despliega, producción, es precisamente el más
importante. Por este motivo, el objetivo
que guía esta sección no es únicamente que el sistema pueda ejecutarse en la
máquina del desarrollador, sino que pueda ejecutarse allí en condiciones lo más
próximas posible a aquellas en las que funcionará una vez desplegado.

\subsection*{Desarrollo en contenedores}
\addcontentsline{toc}{subsection}{Desarrollo en contenedores}

La dificultad de fondo a la hora de reproducir un entorno de ejecución no reside
en el código, sino en todo aquello que lo rodea. Boettiger
\cite{Boettiger2015}, al analizar los obstáculos que impiden reproducir
resultados computacionales, lo expresa con precisión: "instalar o construir el
software necesario para ejecutar el código en cuestión presupone la capacidad de
recrear el entorno computacional de los investigadores originales". La
observación procede del ámbito de la investigación científica, pero describe
igualmente la situación de cualquier equipo de desarrollo: unas instrucciones de
instalación presuponen un conjunto de versiones, dependencias y ajustes del
sistema operativo que rara vez coinciden con los de la máquina en la que se
leen, y ese desajuste se manifiesta como un fallo que solo se produce en algunos
equipos.

El desarrollo en contenedores aborda ese problema desplazando la definición del
entorno del ámbito de la documentación al del código. Según el Instituto
Nacional de Estándares y Tecnología \cite{Souppaya2017}, las tecnologías de
virtualización de aplicaciones se orientan fundamentalmente a proporcionar "una
forma portable, reutilizable y automatizable de empaquetar y ejecutar
aplicaciones", y el término contenedor designa precisamente a estas tecnologías.
La misma publicación señala que su naturaleza portable y declarativa permite a
una organización mantener una gran consistencia entre los entornos de
desarrollo, de pruebas y de producción, hasta el punto de que una imagen creada
en un entorno de desarrollo puede trasladarse a uno de pruebas y copiarse
después en producción sin necesidad de modificación alguna. Aplicado al entorno
local, este enfoque implica que el desarrollador no reconstruye el entorno en su
equipo, sino que instancia el mismo que se emplea en el resto de las etapas de
la entrega.

Entre los mecanismos que implementan este modelo, \textbf{Docker} es el de uso
más extendido y el adoptado en este proyecto. Su documentación oficial lo define
como una plataforma abierta para el desarrollo, la distribución y la ejecución de
aplicaciones, y distingue dos conceptos complementarios: la imagen, que describe
el contenido y la configuración del entorno, y el contenedor, que es una
instancia ejecutable de una imagen.\footnote{Documentación oficial:
	\url{https://docs.docker.com/get-started/docker-overview/}} Esa misma
documentación resume la propiedad que resulta determinante para el entorno
local: los contenedores contienen todo lo necesario para ejecutar la aplicación,
de modo que no es preciso depender de lo que esté instalado en la máquina
anfitriona, y quien recibe un contenedor obtiene exactamente el mismo que
funciona del mismo modo. Es esta garantía, y no la mera comodidad de la
instalación, la que justifica su adopción como base del entorno de desarrollo.

\subsection*{Partes del entorno local}
\addcontentsline{toc}{subsection}{Partes del entorno local}

En el proyecto \textbf{Muncher}, el entorno local no se entiende únicamente como
una réplica de la aplicación desplegada, sino como la suma de dos conjuntos de
elementos: los que componen el sistema en ejecución y los que verifican que un
cambio es correcto. Ambos deben poder levantarse en la máquina del
desarrollador, y ambos deben corresponderse con lo que existe en la nube.

El primer conjunto reproduce la aplicación completa. No se limita al componente
sobre el que se está trabajando, sino que abarca todas las piezas necesarias
para que el sistema funcione de extremo a extremo: el front-end, la API que
atiende sus peticiones y la base de datos que sostiene la información. Cada una
de ellas se corresponde con un servicio de los descritos en
\hyperref[sec:infraestructura]{Infraestructura en la nube}, de manera que la
definición del entorno local no constituye una configuración paralela, sino la
contrapartida local de esa misma arquitectura. Levantar el sistema por partes
—únicamente el front-end, o únicamente la API— resulta insuficiente, ya que
buena parte de los fallos surgen precisamente en la interacción entre
componentes: en la forma de las peticiones, en el esquema de los datos o en la
configuración que los conecta.

El segundo conjunto está formado por los controles de verificación. Un cambio no
se considera terminado porque el sistema arranque, sino cuando supera las
comprobaciones que la integración continua aplicará después: las pruebas
estáticas de formato y análisis del código, las pruebas automatizadas, la
verificación ortográfica y gramatical de los textos de la interfaz y la
construcción de los artefactos desplegables. El requisito determinante es que
estas comprobaciones sean exactamente las mismas en ambos lugares, y no dos
implementaciones equivalentes mantenidas por separado. Cuando se definen una
sola vez y se ejecutan tanto en el equipo del desarrollador como en el
servidor, el resultado obtenido en local anticipa el de la integración
continua; en caso contrario, el desarrollador descubre en el servidor
discrepancias que su propia máquina no era capaz de mostrarle, que es
exactamente la situación que el entorno local pretende evitar.

Ambos conjuntos se ejecutan íntegramente dentro de contenedores Docker. Cada
componente de la aplicación se levanta en el suyo, con la versión de su
\textit{runtime} y sus dependencias fijadas en la imagen correspondiente, y los
controles de verificación se ejecutan asimismo en una imagen que reúne las
herramientas necesarias para aplicarlos. Los únicos requisitos en la máquina del
desarrollador son, por tanto, un sistema de control de versiones y un motor de
contenedores: ni los \textit{runtimes} de los lenguajes empleados, ni los
gestores de dependencias, ni las herramientas de análisis o de formateo se
instalan en el sistema anfitrión.

Esta decisión es la que hace efectivas las dos exigencias planteadas más arriba.
Por un lado, las imágenes que definen los componentes son las mismas que se
despliegan en la nube, de modo que la correspondencia entre el entorno local y
la infraestructura descrita en
\hyperref[sec:infraestructura]{Infraestructura en la nube} deja de depender de
que ambas configuraciones se mantengan sincronizadas a mano. Por otro, la
integración continua no instala herramienta alguna por su cuenta, sino que
construye y utiliza esa misma imagen de verificación, con lo que la identidad
entre los controles locales y los remotos queda garantizada por construcción y
no por convención. El contenedor deja así de ser un mecanismo de empaquetado
para el despliegue y se convierte en el medio ordinario de trabajo durante todo
el desarrollo.
